\section{具体的な活用}
本アプリケーションの実践的な有用性を調査する目的として, プロのストーリークリエーター数グループに, 商業誌に掲載される『ブラック・ジャック』の新作\footnote{TEZUKA2023}のプロットを, 本アプリケーションを用いて制作してもらった.
%本研究では, 開発したアプリケーションの実践的な有用性を調査する目的として, プロのストーリークリエーターに, 出版誌に掲載する『ブラック・ジャック』の新作[TEZUKA 2023]のプロットを, 本アプリケーションを用いて制作してもらう実験を実施した.

プロットの制作期間は, 2023年07月08日から2023年08月02日までのおよそ1ヶ月間.

参加した被験者グループは5グループで, グループAは30代女性1名, グループBは60代男性1名, グループCは40代女性1名と60代男性1名, グループDは60代男性1名, グループEは女性1名と60代男性1名という構成であった.

それぞれのグループにおける創作活動歴と ChatGPT\cite{ref:OpenAI23}等, LLM の使用頻度は表\ref{group}に示す通り. 各被験者グループ内で最も大きな値を示している.

\begin{table}[hbtp]
\begin{center}
{\caption{被験者グループの構成}
\label{group}
\footnotesize
\begin{tabular}{ccc} \hline
グループ & 創作活動歴 & LLMの使用頻度  \\ \hline
A & 14年 & 月に数回使用する\\
B & 37年 & 全く使用したことがない\\
C & 24年 & ほとんど使用しない\\
D & 42年 & 全く使用したことがない\\
E & - & 全く使用したことがない \\
\hline
\end{tabular}
}
\end{center}
\end{table}
\

本実験では, 各被験者グループのプロット制作活動におけるアプリケーションの操作ログを取得した. さらに, 生成されたプロットの一貫性やアプリケーションの操作性, アプリケーションを活用したことによる創造性への刺激などについて, 5段階のリッカート尺度および自由記述により評価を求めるアンケートを実施した.